%% =============================================================================
%%  reproduction.tex -- interim results report, IEEE double column.
%%
%%  Build:  pdflatex reproduction.tex   (twice, for cross-references)
%%
%%  This is a working report of results obtained so far, not a submission draft.
%%  Sections marked "in progress" are exactly that, and the numbers behind every
%%  table are regenerated by a committed script named in the caption.
%% =============================================================================
\documentclass[conference]{IEEEtran}
\IEEEoverridecommandlockouts

\usepackage{cite}
\usepackage{amsmath,amssymb}
\usepackage{graphicx}
\usepackage{booktabs}
\usepackage{url}
\usepackage[hidelinks]{hyperref}
\usepackage{xcolor}

% Marks a claim whose evidence is still being collected, so a reader is never
% left guessing which numbers are final.
\newcommand{\pending}[1]{\textcolor{red!70!black}{#1}}

\begin{document}

\title{Capacity, Not Associativity: A Reproduction Study of a
Seamless Data-Cache Extension for Multi-Core RISC-V}

\author{\IEEEauthorblockN{Krishna Gorai}
\IEEEauthorblockA{Indian Institute of Technology Mandi\\
Mandi, Himachal Pradesh, India}
\and
\IEEEauthorblockN{Bikram Paul}
\IEEEauthorblockA{Indian Institute of Technology Mandi\\
Mandi, Himachal Pradesh, India}}

\maketitle

\begin{abstract}
We reproduce, from scratch and in open-source SystemVerilog, the snoop-coherent
private L1 data-cache extension proposed by Kamaleldin \emph{et al.} for
FPGA-based multi-core RISC-V systems. The reproduction integrates an unmodified
CV32E40P core, a write-through no-allocate data-cache unit, a snooping
invalidation bus and an AXI4 crossbar into a four-processing-element system, and
implements it on a Zynq UltraScale+ ZCU104. Reproducing the original's hit-rate
figure required no change to the hardware but a specific change to how the
benchmark kernels write their results, which we report as a reproducibility
finding. Our principal result concerns the original's FFT curve: a controlled
sweep that holds capacity fixed while doubling associativity, and then the
reverse, shows the reported flat curve does not follow from the stated
2\,KiB cache. Capacity recovers the curve completely while associativity
recovers a sixth of it at $N{=}512$ and almost nothing at $N{=}1024$, a ratio of
$6.6{:}1$ widening to $44{:}1$. We further show the extension's benefit is contingent on reuse and
memory latency rather than intrinsic: on a streaming kernel at low memory
latency the extension is slower than no cache at all, and the same kernel
reaches $2.285\times$ when only the latency is raised. The complete design,
data and regeneration scripts are public.
\end{abstract}

\begin{IEEEkeywords}
RISC-V, cache coherence, snooping protocols, FPGA, reproducibility
\end{IEEEkeywords}

% =============================================================================
\section{Introduction}
% =============================================================================

Attaching a data cache to a soft RISC-V core without modifying the core is an
attractive proposition for FPGA-based multi-core systems: the core remains a
verified, unmodified IP block, and coherence is handled entirely in hardware.
Kamaleldin \emph{et al.}~\cite{kamaleldin2024} propose exactly this, a private
L1 data-cache sub-system attached to an in-order RISC-V core through its
existing memory interface, with a snooping invalidation bus maintaining
coherence and no software involvement.

We set out to reproduce that design in open-source RTL. This report presents the
results obtained so far. Our contributions are:

\begin{itemize}
  \item A from-scratch, open-source SystemVerilog implementation of the
        sub-system, integrated into a four-PE RISC-V system and implemented on a
        ZCU104, with the non-coherent baseline built from the same sources so
        the difference between them is attributable.
  \item A controlled measurement showing that the flat FFT hit-rate curve
        reported in the original does not follow from the cache configuration
        the original states, and that the effect is one of capacity rather than
        conflict (Section~\ref{sec:capacity}).
  \item A measurement bounding the extension's benefit: it is contingent on
        reuse and memory latency, and is negative on a streaming kernel at low
        latency (Section~\ref{sec:memlat}).
  \item A reproducibility finding: the original's execution-time figure cannot
        be reproduced with kernels that accumulate results in a register, for a
        reason specific to write-through no-allocate caches, and requires no
        hardware change to resolve (Section~\ref{sec:fig6}).
  \item Six protocol and implementation hazards encountered in building the
        design that the original's description does not cover
        (Section~\ref{sec:bugs}).
\end{itemize}

% =============================================================================
\section{The Design Under Reproduction}
% =============================================================================

The original attaches a Data Cache Unit (DCU) between an in-order RISC-V core
and the memory system. The DCU is a private, write-through, no-allocate L1 data
cache. Coherence is maintained by a snooping invalidation bus: a write broadcasts
an invalidation to the other DCUs, which drop any matching line. Because the
policy is write-through, memory always holds the current value and no dirty-line
writeback or ownership state is required. The core is unmodified and unaware of
the cache, which is the origin of the ``seamless'' claim.

Our reproduction targets the stated configuration: a 2\,KiB DCU, two-way set
associative, 64 sets of 16-byte lines, one per processing element, with four PEs
in the system.

% =============================================================================
\section{Implementation}
% =============================================================================

Each processing element comprises an unmodified CV32E40P core (pinned as a git
submodule, with \texttt{COREV\_PULP}, \texttt{COREV\_CLUSTER} and the FPU
disabled), a private instruction tightly-coupled memory, an
Instruction-Bridge and a Data-Bridge. The core issues ordinary loads and stores;
address decoding in the Data-Bridge routes coherent traffic to the DCU and
everything else past it.

The four PEs meet the shared memories through a real AXI4 crossbar rather than a
point-to-point arbiter, because the original lists the interconnect as a
separate resource line and a simplified fabric would not be comparable. Each PE
presents three masters onto three slaves. AXI4 rather than AXI4-Lite is
deliberate: a line fill is one INCR burst of four beats, so a miss pays one
arbitration and one address phase rather than four.

A non-coherent baseline is built from the same sources with a single parameter
changed, substituting a bypass shim for each DCU. Every comparison in this
report is between two builds that differ in that one parameter and nothing else.

% =============================================================================
\section{Verification}
% =============================================================================

The design is verified by self-checking testbenches at three levels: unit
(the DCU and the AXI fabric in isolation), subsystem (four DCUs and the snooping
bus against a version-monotonicity coherence checker), and system (the full
four-PE design running bare-metal kernels). At the system level the coherent and
non-coherent builds run identical binaries and are required to produce identical
checksums, so any coherence failure appears as a value mismatch rather than as a
performance difference.

% =============================================================================
\section{Results}
% =============================================================================

\subsection{Functional reproduction}

The four-PE system passes 273 self-checking assertions on a parallel
bring-up kernel in both the coherent and non-coherent configurations, with
identical results. Cache statistics are reported in the split the original's
hit-rate figure uses: on this kernel the coherent build records a
85.35\,\% read hit rate and a 93.77\,\% write hit rate.

\subsection{Reproducing the execution-time figure}
\label{sec:fig6}

The original reports execution-time reductions of approximately 40, 40 and
23\,\% for its three kernels. Our first attempts fell far short, and the cause
was neither the cache nor the baseline, both of which were already correct.

Under a write-through no-allocate policy, a store to a line not currently
resident neither allocates nor hits. A kernel that accumulates its result in a
register and writes it once at the end therefore produces a write hit rate
pinned near zero regardless of how well the cache is working, because the single
write is always to a line the cache has no reason to hold. The kernels must
accumulate \emph{in memory}, and the matrix multiplication must use i-k-j
ordering so that the accumulating row stays resident.

With that change and appropriately sized working sets, and \emph{no change to
the RTL}, we obtain reductions of 38.4, 48.0 and 18.5\,\% against the original's
40, 40 and 23\,\%. We report this because it is invisible from the original's
description: the figure is reproducible, but only with a benchmark writing
discipline the paper does not state.

\subsection{Capacity versus associativity}
\label{sec:capacity}

The original's FFT curve is reported as flat across problem sizes at the stated
2\,KiB DCU. We do not observe a flat curve.

The per-PE FFT working set is $8N$ bytes, so $N{=}512$ is the first size that
exceeds a 2\,KiB cache. It is also the first size at which the butterfly
half-span reaches the set-index period, so the two operands of a butterfly alias
onto one set. Capacity pressure and conflict aliasing arrive together and a
single run cannot distinguish them. Four configurations can: holding capacity at
2\,KiB while doubling associativity isolates conflict, and holding associativity
at two ways while quadrupling capacity isolates capacity.

\begin{table}[t]
\caption{Read hit rate, FFT, coherent build. All four configurations are
bit-identical at the two control sizes, so only the parameter under test moves.
Regenerated by \texttt{scripts/sweep\_capacity.sh}.}
\label{tab:capacity}
\centering
\begin{tabular}{@{}rrrrrr@{}}
\toprule
$N$ & Working set & 2-way & 4-way & 2-way & 2-way \\
    &             & 2\,KiB & 2\,KiB & 8\,KiB & 16\,KiB \\
\midrule
128  & 1\,KiB & 96.8 & 96.8 & 96.8 & 96.8 \\
256  & 2\,KiB & 97.2 & 97.2 & 97.2 & 97.2 \\
512  & 4\,KiB & \textbf{66.0} & 70.8 & \textbf{97.5} & 97.5 \\
1024 & 8\,KiB & \textbf{57.8} & 58.7 & \textbf{97.7} & 97.7 \\
\bottomrule
\end{tabular}
\end{table}

Table~\ref{tab:capacity} gives the result. At the two sizes whose working set
fits, all four configurations are bit-identical, which establishes that nothing
but the parameter under test is moving. At $N{=}512$ the stated configuration
loses 31.2 points of read hit rate. Doubling associativity returns 4.8 of them;
quadrupling capacity returns 31.5, restoring the pre-overflow rate. At $N{=}1024$
associativity returns 0.9 points and capacity returns 39.9.

Capacity therefore outweighs associativity by $6.6{:}1$ at $N{=}512$ and
$44{:}1$ at $N{=}1024$. The ratio widens with the working set because additional
ways cannot help once nothing fits. In execution time the stated configuration
costs 24.5\,\% at $N{=}512$ and 27.6\,\% at $N{=}1024$ relative to a cache that
holds the working set.

The observation that makes this a capacity argument rather than an arbitrary one
is the final column: 16\,KiB is bit-identical to 8\,KiB at every size. Capacity
helps until the working set fits and then stops entirely. A conflict-dominated
effect would not saturate in that way.

We therefore report that the flat FFT curve does not follow from a 2\,KiB DCU. A
cache that holds the $8N$-byte per-PE working set produces it; the stated
configuration produces a 31-point collapse at $N{=}512$ and a 39-point collapse
at $N{=}1024$.

\subsection{What the extension is actually worth}
\label{sec:memlat}

The reported benefit of a cache extension is not a property of the extension
alone. The DCU introduces a second pipeline stage, so every access costs one
cycle it did not cost before, plus a snoop arbitration; the non-coherent path
pays neither. That overhead is fixed. The benefit, by contrast, has two terms:
how often a line is touched again before eviction, and how much memory latency
each of those hits avoids. Neither term alone predicts the outcome.

\begin{table}[t]
\caption{Speedup of the coherent build over the non-coherent baseline, against
modelled shared-memory latency. Regenerated by
\texttt{scripts/sweep\_memlat.sh}.}
\label{tab:memlat}
\centering
\begin{tabular}{@{}lrrr@{}}
\toprule
Kernel & MemLat 2 & MemLat 8 & MemLat 20 \\
\midrule
memcpy (16\,KiB)     & \textbf{0.975} & 1.509 & 2.285 \\
FFT ($N{=}256$)      & 1.329 & 2.267 & 4.425 \\
matmul ($64^2$)      & 1.638 & 3.067 & 4.570 \\
conv2d ($64^2$)      & 1.852 & 4.133 & \textbf{7.941} \\
\bottomrule
\end{tabular}
\end{table}

Table~\ref{tab:memlat} sweeps both builds over three memory latencies. Speedup
is monotonic in latency for every kernel, and the ordering between kernels is
preserved at all three latencies, which indicates the two terms act
independently rather than interacting.

The first row is the important one. At a memory latency of two cycles the
extension makes memcpy \emph{slower}: 21\,010 cycles against the baseline's
20\,493, a speedup of $0.975\times$. The same kernel, unchanged, reaches
$2.285\times$ at a latency of twenty. Reuse is constant across those three
runs, so latency alone moves the result across the break-even line. A
cross-kernel comparison cannot establish that, because reuse and latency vary
together; a single kernel crossing over does.

The mechanism is visible in the cycle counts rather than the ratios. Raising
latency from 2 to 20 costs the non-coherent conv2d build 417\,\% more cycles
and the coherent build 21\,\%; for FFT the figures are 269\,\% and 11\,\%. The
cached builds barely register the change. What the extension does is convert a
latency-bound workload into a compute-bound one, and the speedup therefore
measures chiefly how latency-bound the workload was to begin with. Where there
is no reuse and no latency to hide, as in memcpy at MemLat 2, there is nothing
to convert and the overhead is all that remains.

We note also that the growth is not unbounded. matmul rises $1.638 \rightarrow
3.067 \rightarrow 4.570$, decelerating as the cached build's residual misses
come to dominate its own runtime, at which point further latency penalises both
builds proportionally.

This bounds the original's headline result rather than contradicting it. The
reductions the original reports are reproducible, but they are contingent on the
memory latency of the platform, and a system with fast shared memory and a
streaming access pattern can lose from the extension outright.

\subsection{FPGA implementation}

Both configurations were synthesised and implemented on a Zynq UltraScale+
XCZU7EV (ZCU104) in a non-project Vivado 2025.1 flow, and independently through
a managed project, at a 75\,MHz system clock derived from the 300\,MHz board
clock by an integer divider.

\begin{table}[t]
\caption{Post-implementation resources, coherent build, XCZU7EV at 75\,MHz.
Regenerated by \texttt{scripts/table1.py}.}
\label{tab:resources}
\centering
\begin{tabular}{@{}lrrrr@{}}
\toprule
Module & $\times$ & LUTs & FFs & BRAM \\
\midrule
Processing element        & 4  & 17\,837 & 9\,948 & 32 \\
\quad RISC-V core         & 4  & 16\,937 & 8\,404 & 0 \\
\quad ITCM + I-Bridge     & 4  & 253     & 36     & 32 \\
\quad Data-Bridge         & 4  & 319     & 32     & 0 \\
\quad AXI4 masters        & 12 & 328     & 1\,476 & 0 \\
\addlinespace
\textbf{Cache extension}  & 1  & \textbf{5\,322} & \textbf{4\,118} & \textbf{20} \\
\quad Data Cache Unit     & 4  & 5\,019  & 4\,000 & 20 \\
\quad\quad controller     & 4  & 1\,697  & 2\,848 & 0 \\
\quad\quad tag RAM        & 4  & 975     & 96     & 4 \\
\quad\quad data RAM       & 4  & 2\,364  & 1\,056 & 16 \\
\quad Snoopy bus          & 1  & 303     & 118    & 0 \\
\addlinespace
AXI4 crossbar             & 1  & 2\,127  & 279    & 0 \\
Shared data memory        & 1  & 724     & 147    & 64 \\
Shared instruction memory & 1  & 506     & 86     & 8 \\
Control region            & 1  & 114     & 259    & 0 \\
\midrule
Total system              &    & 26\,635 & 15\,155 & 124 \\
\bottomrule
\end{tabular}
\end{table}

Table~\ref{tab:resources} gives the per-module breakdown. The cache extension
costs 5\,117 LUTs and 3\,487 flip-flops over the non-coherent baseline
(26\,635 against 21\,518 LUTs, 15\,155 against 11\,668 flip-flops), together
with 20 additional block RAM tiles and no additional DSPs. Two observations are
worth drawing out. First, the extension is 20\,\% of system LUTs against the
four cores' 64\,\%, so it is not the dominant area cost. Second, within the DCU
the controller outweighs the storage, 1\,697 LUTs of control against 975 of tag
RAM, because at 2\,KiB the arrays are nearly free and the coherence logic is
not.

\begin{table}[t]
\caption{Post-implementation timing, XCZU7EV, 75\,MHz target.}
\label{tab:timing}
\centering
\begin{tabular}{@{}lrr@{}}
\toprule
 & Coherent & Baseline \\
\midrule
Worst negative slack (setup) & $+0.880$\,ns & $+0.929$\,ns \\
Worst hold slack             & $+0.003$\,ns & $+0.010$\,ns \\
Failing endpoints            & 0 & 0 \\
Timing constraints met       & yes & yes \\
\bottomrule
\end{tabular}
\end{table}

Timing is given in Table~\ref{tab:timing}. Both configurations meet every
constraint. Reaching that state required one change to the board wrapper: the
single register driving every asynchronous reset in the system fans out to
approximately fifteen thousand clear pins, and the removal check on reset
release failed on routing distance alone, with 97\,\% of the worst path in wire.
Constraining that register's fan-out so the tool replicates it removes the
violation. We note that before this change the coherent build met timing by
12\,ps in one placement and failed the same check in two others: a result that
was a property of the placement rather than of the design.

\subsection{Post-implementation verification}

The implemented design was verified by simulating the placed-and-routed netlist.
A board-level testbench drives only the differential board clock, the reset
button and the run switch, and loads nothing: the program reaches the design as
block RAM initialisation contents established at synthesis. All four processing
elements complete and report success, and the netlist tracks the register-level
simulation cycle for cycle, confirming that implementation introduced no
functional divergence.

% =============================================================================
\section{Protocol and Implementation Hazards}
\label{sec:bugs}
% =============================================================================

Building the design surfaced hazards that its published description does not
cover. We report them because they are the substance of making such a design
correct.

\begin{enumerate}
  \item \textbf{Shared array read port against a stalled stage.} The tag and
        data array read address is driven by the first pipeline stage; while the
        second stalled on a miss, the first stage's differing index displaced the
        data the second was still evaluating. Every hit decision taken after a
        stall used the wrong set.
  \item \textbf{Coherence lock released before the write lands.} Releasing the
        single-writer lock when the invalidation is \emph{granted} is
        insufficient. A remote read granted between that point and the
        write-through actually reaching memory fetches the stale value and caches
        it, and no further invalidation is ever sent. The lock must be held until
        the write lands.
  \item \textbf{A failed store-conditional releasing another core's
        reservation.} With one link register per cluster, a store-conditional may
        only consume the reservation if it owns it.
  \item \textbf{The non-coherent baseline never withdrawing a granted request.}
        A request must be dropped on grant, not on response. The cached path had
        always done so; the baseline had not, and surplus reads returned data
        into whatever transaction was outstanding by then.
  \item \textbf{Combinational blocks re-triggering on their own writes.} A
        bit-write through a variable index is a read-modify-write of the whole
        vector, which places the vector in its own block's sensitivity list. The
        block reschedules forever and simulation time stops, with no error and
        nothing visible in a waveform, because the same values are reassigned.
  \item \textbf{A helper function reading a port rather than taking it as an
        argument}, so the port never entered the sensitivity list of a continuous
        assignment and every invalidation carried a stale address.
\end{enumerate}

% =============================================================================
\section{Threats to Validity}
% =============================================================================

Our claim in Section~\ref{sec:capacity} is that a published result does not
follow from its stated configuration, and it rests on our reading of that
configuration. A different line size, a shared rather than private cache
organisation, or a different definition of the FFT working set would each
explain the discrepancy without any error in the original. We state our
assumptions explicitly for that reason: 2\,KiB per PE, private, two-way, 16-byte
lines, an $8N$-byte per-PE working set for a complex single-precision FFT of
size $N$.

The original targets a Virtex UltraScale+ XCVU9P at 60\,MHz and this work
targets an XCZU7EV at 75\,MHz, so absolute resource and frequency figures are
directional rather than a like-for-like comparison. Our processing element
presents three AXI masters where the original's figure implies two, which
inflates our crossbar's resource line relative to theirs.

Power is not reported. Vivado's vectorless estimate models this design as idle,
because it cannot evaluate the clock-gate enables on the cores, and a
switching-activity-based estimate has not yet been completed.

% =============================================================================
\section{Conclusion}
% =============================================================================

We have reproduced a seamless snoop-coherent data-cache extension for multi-core
RISC-V in open-source RTL and implemented it on an FPGA, and in doing so found
that the original's flat FFT hit-rate curve does not follow from its stated
2\,KiB cache. The effect is one of capacity, not conflict: capacity outweighs
associativity by $6.6{:}1$ and then $44{:}1$ as the working set grows, and
additional capacity stops helping precisely when the working set fits. We further show that the extension's benefit is contingent on reuse and memory
latency: it converts a latency-bound workload into a compute-bound one, and
where there is neither reuse nor latency to hide it is slower than no cache at
all. The design, data and regeneration scripts are available at
\url{https://github.com/Krishna-Gorai/riscv-cache-extension}.

\begin{thebibliography}{1}
\bibitem{kamaleldin2024}
A.~Kamaleldin, M.~Nickel, S.~Wu and D.~G\"ohringer,
``Seamless Cache Extension for FPGA-based Multi-Core RISC-V SoC,''
in \emph{Proc. IEEE 37th Int. System-on-Chip Conf. (SOCC)}, 2024,
doi:10.1109/SOCC62300.2024.10737850.
\end{thebibliography}

\end{document}
